\chapter{Chapter07 — Complexity}
%%% generated by the blueprint pipeline from FABL.Chapter07.Complexity
\section{Overview}

Declarations of the Lean module \texttt{FABL.Chapter07.Complexity}.

\section{Declarations}

\begin{definition}[Polynomial-time map]
\lean{FABL.PolytimeMap}
\statementsource{BGS75}{pdf-2459e3840724-p003-b004}
\label{FABL.PolytimeMap}
\leanok
A polynomial-time map on binary strings consists of a function from finite binary
strings to finite binary strings together with a certificate that a single-tape Turing
machine computes that function in polynomial time.
\end{definition}

\begin{theorem}[Monotonicity of polynomial evaluation]
\lean{FABL.PolytimeMap.polynomial_eval_monotone}
\label{FABL.PolytimeMap.polynomial_eval_monotone}
\leanok
\difficulty{3}
Every polynomial with natural-number coefficients is monotone on the natural numbers.
\end{theorem}

\begin{definition}[Polynomial-Time Certificate Normalization]
\lean{FABL.PolytimeMap.monotoneBounded}
\label{FABL.PolytimeMap.monotoneBounded}
\leanok
For a function $f$ from Boolean lists to Boolean lists and a certificate that $f$ is
computable in polynomial time, the certificate obtained by taking the same machine and
polynomial uses that polynomial, evaluated at the input length, as its time bound.
\end{definition}

\begin{theorem}[Monotonicity of the runtime bound]
\lean{FABL.PolytimeMap.monotone_monotoneBounded}
\label{FABL.PolytimeMap.monotone_monotoneBounded}
\leanok
\uses{FABL.PolytimeMap.monotoneBounded, FABL.PolytimeMap.polynomial_eval_monotone}
\difficulty{1}
If a function on Boolean lists is computable in polynomial time, then the polynomial
runtime bound obtained by evaluating its runtime polynomial is monotone as a function of
the input length.
\end{theorem}

\begin{definition}[Identity polynomial-time map]
\lean{FABL.PolytimeMap.id}
\label{FABL.PolytimeMap.id}
\leanok
\uses{FABL.PolytimeMap}
The identity function on binary strings, which returns each input string unchanged, is
equipped with its polynomial-time computation certificate.
\end{definition}

\begin{definition}[Composition of polynomial-time maps]
\lean{FABL.PolytimeMap.comp}
\label{FABL.PolytimeMap.comp}
\leanok
\uses{FABL.PolytimeMap, FABL.PolytimeMap.monotoneBounded, FABL.PolytimeMap.monotone_monotoneBounded}
Given polynomial-time maps $f$ and $g$ on binary strings, their composition is the map
that first applies $f$ and then applies $g$; it is equipped with a polynomial-time
computation certificate.
\end{definition}

\begin{theorem}[Application of the identity map]
\lean{FABL.PolytimeMap.id_apply}
\label{FABL.PolytimeMap.id_apply}
\leanok
\uses{FABL.PolytimeMap.id}
\difficulty{1}
For every Boolean list $x$, the identity map sends $x$ to $x$.
\end{theorem}

\begin{theorem}[Application of composed maps]
\lean{FABL.PolytimeMap.comp_apply}
\label{FABL.PolytimeMap.comp_apply}
\leanok
\uses{FABL.PolytimeMap, FABL.PolytimeMap.comp}
\difficulty{1}
For polynomial-time maps $F,G$ on Boolean lists and any Boolean list $x$, the
composition $G\circ F$ sends $x$ to $G(F(x))$.
\end{theorem}

\begin{definition}[Canonical binary encoding]
\lean{FABL.BinaryEncoding}
\label{FABL.BinaryEncoding}
\leanok
A binary encoding of a type of objects assigns each object a binary string and provides
a decoder from binary strings to either an object or failure. Decoding an object's
encoding recovers that object, and every successfully decoded string is the encoding of
the decoded object.
\end{definition}

\begin{theorem}[Injectivity of binary encoding]
\lean{FABL.BinaryEncoding.encode_injective}
\label{FABL.BinaryEncoding.encode_injective}
\leanok
\uses{FABL.BinaryEncoding}
\difficulty{3}
Let $A$ be a type with an encoding $e:A\to\{0,1\}^*$ and a partial decoder
$d:\{0,1\}^*\to A$. Suppose decoding an encoded object recovers it, so $d(e(a))=a$ for
every $a\in A$, and every successfully decoded string is the encoding of its decoded
object. Then $e$ is injective.
\end{theorem}

\begin{definition}[Language of valid encodings]
\lean{FABL.BinaryEncoding.language}
\label{FABL.BinaryEncoding.language}
\leanok
\uses{FABL.BinaryEncoding}
Let $A$ be a type with a binary encoding and decoding map, such that decoding an encoded
object recovers it and every successfully decoded string is the encoding of its decoded
object. The language of the encoding consists of exactly those finite binary strings
that decode successfully.
\end{definition}

\begin{theorem}[Encoded strings belong to the decoding language]
\lean{FABL.BinaryEncoding.encode_mem_language}
\label{FABL.BinaryEncoding.encode_mem_language}
\leanok
\uses{FABL.BinaryEncoding, FABL.BinaryEncoding.language}
\difficulty{3}
Let $A$ be a type with an encoding $e:A\to\{0,1\}^*$ and a partial decoder
$d:\{0,1\}^*\to A$. Suppose decoding an encoded object recovers it, so $d(e(a))=a$ for
every $a\in A$, and every successfully decoded string is the encoding of its decoded
object. The language consisting of strings on which $d$ is defined contains $e(a)$ for
every $a\in A$.
\end{theorem}

\begin{definition}[Packed Boolean circuit]
\lean{FABL.PackedBooleanCircuit}
\label{FABL.PackedBooleanCircuit}
\leanok
\uses{FABL.BooleanCircuit}
A packed Boolean circuit consists of a positive number $n$ of input bits and a Boolean
circuit on those $n$ bits. The circuit is an ordered sequence of gates, each of which
reads an input bit, is a Boolean constant, or applies negation, conjunction, or
disjunction to earlier gates; one gate is distinguished as the output.
\end{definition}

\begin{definition}[Packed circuit size]
\lean{FABL.PackedBooleanCircuit.size}
\label{FABL.PackedBooleanCircuit.size}
\leanok
\uses{FABL.BooleanCircuit.size, FABL.PackedBooleanCircuit}
For a packed Boolean circuit with $n$ input bits and $g$ gates, define its size to be
$n+g$.
\end{definition}

\begin{definition}[Boolean circuit encoding]
\lean{FABL.BooleanCircuitEncoding}
\label{FABL.BooleanCircuitEncoding}
\leanok
\uses{FABL.BinaryEncoding, FABL.PackedBooleanCircuit, FABL.PackedBooleanCircuit.size}
A Boolean circuit encoding assigns binary strings to packed Boolean circuits and decodes
every such string back to its circuit, with encoding followed by decoding recovering
each circuit and every successfully decoded string being the canonical encoding of its
decoded circuit. The length of each circuit's encoding is at least its size, where size
is the number of input bits plus the number of gates.
\end{definition}

\begin{definition}[Packed PCPP description]
\lean{FABL.PackedPCPPDescription}
\label{FABL.PackedPCPPDescription}
\leanok
\uses{FABL.BooleanCircuit.acceptedProperty, FABL.PCPPDescription, FABL.PackedBooleanCircuit, FABL.SemanticPCPPReduction}
Let a semantic reduction assign to each Boolean circuit on a positive number $n$ of
input bits a PCPP description for the property that the circuit accepts. A packed PCPP
description consists of a packed Boolean circuit and a PCPP description for its
accepted-input property, with the description equal to the one assigned to that circuit
by the semantic reduction.
\end{definition}

\begin{definition}[Packing a PCPP description]
\lean{FABL.packPCPPDescription}
\label{FABL.packPCPPDescription}
\leanok
\uses{FABL.PackedBooleanCircuit, FABL.PackedPCPPDescription, FABL.SemanticPCPPReduction}
Let a semantic reduction assign to every Boolean circuit with positive input length a
PCPP description for the property accepted by that circuit, and let $C$ be a Boolean
circuit with positive input length. The packed PCPP description built from $C$ has
source circuit $C$ and description equal to the one assigned to $C$ by the reduction.
\end{definition}

\begin{definition}[Two-parameter polynomial time bound]
\lean{FABL.PolynomialTimeBound₂}
\leanok
A two-parameter polynomial time bound consists of a function
$T\colon\mathbb{N}^2\to\mathbb{N}$ and a polynomial $p$ with natural-number coefficients
such that, for all $a,b\in\mathbb{N}$, $T(a,b)\le p(a+b)$.
\end{definition}

\begin{definition}[Output-sensitive PCPP reduction]
\lean{FABL.OutputSensitivePCPPReduction}
\statementsource{Analysis-of-Boolean-Functions-by-Ryan-ODonnell}{
  pdf-ffbe9c7b3306-p182-b004,
  pdf-ffbe9c7b3306-p183-b001
}
\label{FABL.OutputSensitivePCPPReduction}
\leanok
\uses{FABL.BinaryEncoding, FABL.BooleanCircuitEncoding, FABL.PackedBooleanCircuit.size, FABL.PackedPCPPDescription, FABL.PolynomialTimeBound₂, FABL.SemanticPCPPReduction, FABL.packPCPPDescription}
Fix a canonical binary encoding of positive-input-length Boolean circuits whose encoding
length is at least the circuit size, where size is the number of input bits plus the
number of gates. An output-sensitive PCPP reduction consists of a semantic assignment of
a PCPP description to each such circuit, a canonical binary encoding of the resulting
packed descriptions, and a total function on binary strings that agrees with this
encoding on encoded circuits. It also includes a single-tape Turing machine and a
two-parameter runtime bound majorized by a polynomial in the sum of its arguments; on
every input the machine produces the function’s output within that bound, using the
circuit size and its generated proof length when decoding succeeds, and the input length
and zero when decoding fails.
\end{definition}

\begin{definition}[Polynomial bound on proof length]
\lean{FABL.ProofLengthPolynomiallyBounded}
\statementsource{Analysis-of-Boolean-Functions-by-Ryan-ODonnell}{pdf-ffbe9c7b3306-p183-b004}
\label{FABL.ProofLengthPolynomiallyBounded}
\leanok
\uses{FABL.BooleanCircuitEncoding, FABL.OutputSensitivePCPPReduction, FABL.PackedBooleanCircuit, FABL.PackedBooleanCircuit.size}
Fix a Boolean circuit encoding and an output-sensitive PCPP reduction for that encoding.
There is a polynomial with natural-number coefficients such that, for every packed
Boolean circuit, the proof length of the reduction’s PCPP description is at most the
polynomial evaluated at the circuit’s size.
\end{definition}

\begin{definition}[Polynomial-time computability of the reduction output]
\lean{FABL.OutputSensitivePCPPReduction.toPolyTimeComputable}
\statementsource{Analysis-of-Boolean-Functions-by-Ryan-ODonnell}{
  pdf-ffbe9c7b3306-p182-b004,
  pdf-ffbe9c7b3306-p183-b001
}
\label{FABL.OutputSensitivePCPPReduction.toPolyTimeComputable}
\leanok
\uses{FABL.BooleanCircuitEncoding, FABL.OutputSensitivePCPPReduction, FABL.PackedBooleanCircuit.size, FABL.PolytimeMap.polynomial_eval_monotone, FABL.ProofLengthPolynomiallyBounded}
Fix a Boolean circuit encoding and an output-sensitive PCPP reduction, and suppose its
generated proof lengths are bounded by a polynomial in circuit size. Then the
reduction’s total output function on binary strings is computed by a single-tape Turing
machine in polynomial time.
\end{definition}

\begin{definition}[Polynomial-time map from the reduction]
\lean{FABL.OutputSensitivePCPPReduction.toPolytimeMap}
\statementsource{Analysis-of-Boolean-Functions-by-Ryan-ODonnell}{
  pdf-ffbe9c7b3306-p182-b004,
  pdf-ffbe9c7b3306-p183-b001
}
\label{FABL.OutputSensitivePCPPReduction.toPolytimeMap}
\leanok
\uses{FABL.BooleanCircuitEncoding, FABL.OutputSensitivePCPPReduction, FABL.OutputSensitivePCPPReduction.toPolyTimeComputable, FABL.PolytimeMap, FABL.ProofLengthPolynomiallyBounded}
Fix a Boolean circuit encoding and an output-sensitive PCPP reduction, and suppose its
generated proof lengths are bounded by a polynomial in circuit size. The reduction’s
output function, viewed as a map from finite binary strings to finite binary strings, is
a polynomial-time map.
\end{definition}

\begin{definition}[Polynomial-time many-one reducibility]
\lean{FABL.KarpReducible}
\statementsource{Yap83}{
  pdf-1a201b302069-p001-b004,
  pdf-1a201b302069-p003-b001
}
\statementsource{BGS75}{pdf-2459e3840724-p003-b004}
\label{FABL.KarpReducible}
\leanok
\uses{FABL.PolytimeMap}
A language $L_1$ is polynomial-time many-one reducible to a language $L_2$ if there is a
polynomial-time map $f$ on binary strings such that, for every binary string $x$, $x\in
L_1$ if and only if $f(x)\in L_2$.
\end{definition}

\begin{theorem}[Reflexivity of Karp reducibility]
\lean{FABL.KarpReducible.refl}
\label{FABL.KarpReducible.refl}
\leanok
\uses{FABL.PolytimeMap.id}
\difficulty{1}
Every Boolean language is Karp reducible to itself.
\end{theorem}

\begin{theorem}[Transitivity of Karp reducibility]
\lean{FABL.KarpReducible.trans}
\label{FABL.KarpReducible.trans}
\leanok
\uses{FABL.PolytimeMap.comp}
\difficulty{4}
If three languages over the Boolean alphabet satisfy that the first is Karp reducible to
the second and the second is Karp reducible to the third, then the first is Karp
reducible to the third.
\end{theorem}

\begin{definition}[Binary promise problem]
\lean{FABL.BinaryPromiseProblem}
\label{FABL.BinaryPromiseProblem}
\leanok
A binary promise problem consists of a set of Boolean strings designated as
yes-instances and a set designated as no-instances, with the two sets disjoint. Strings
outside these sets need not have either designation.
\end{definition}

\begin{definition}[Promise-preserving polynomial reduction]
\lean{FABL.PromiseKarpReducible}
\label{FABL.PromiseKarpReducible}
\leanok
\uses{FABL.BinaryPromiseProblem, FABL.PolytimeMap}
For binary promise problems $A$ and $B$, a polynomial-time reduction from $A$ to $B$ is
a polynomial-time computable map on Boolean strings that sends every yes-instance of $A$
to a yes-instance of $B$ and every no-instance of $A$ to a no-instance of $B$.
\end{definition}

\begin{theorem}[Reflexivity of promise reductions]
\lean{FABL.PromiseKarpReducible.refl}
\label{FABL.PromiseKarpReducible.refl}
\leanok
\uses{FABL.BinaryPromiseProblem, FABL.PolytimeMap.id, FABL.PromiseKarpReducible}
\difficulty{1}
Every binary promise problem reduces to itself by a polynomial-time map that preserves
its promised yes-instances and no-instances.
\end{theorem}

\begin{theorem}[Transitivity of promise reductions]
\lean{FABL.PromiseKarpReducible.trans}
\label{FABL.PromiseKarpReducible.trans}
\leanok
\uses{FABL.BinaryPromiseProblem, FABL.PolytimeMap.comp, FABL.PromiseKarpReducible}
\difficulty{4}
Let $A$, $B$, and $C$ be binary promise problems. If $A$ reduces to $B$ and $B$ reduces
to $C$ by polynomial-time maps preserving promised yes-instances and no-instances, then
$A$ reduces to $C$ by such a map.
\end{theorem}

\begin{definition}[Polynomial-time promise distinguisher]
\lean{FABL.PolytimeMap.DistinguishesPromise}
\label{FABL.PolytimeMap.DistinguishesPromise}
\leanok
\uses{FABL.BinaryPromiseProblem, FABL.PolytimeMap}
A polynomial-time map distinguishes a binary promise problem if it outputs the one-bit
string $[\mathrm{true}]$ on every yes-instance and the one-bit string $[\mathrm{false}]$
on every no-instance. There is no condition on its outputs outside the promised sets.
\end{definition}

\begin{definition}[Existence of a polynomial-time distinguisher]
\lean{FABL.HasPolytimeDistinguisher}
\label{FABL.HasPolytimeDistinguisher}
\leanok
\uses{FABL.BinaryPromiseProblem, FABL.PolytimeMap, FABL.PolytimeMap.DistinguishesPromise}
A binary promise problem has a polynomial-time distinguisher if some polynomial-time map
outputs the one-bit string $[\mathrm{true}]$ on every yes-instance and
$[\mathrm{false}]$ on every no-instance.
\end{definition}

\begin{theorem}[Distinguisher transfer under reduction]
\lean{FABL.PromiseKarpReducible.hasPolytimeDistinguisher}
\label{FABL.PromiseKarpReducible.hasPolytimeDistinguisher}
\leanok
\uses{FABL.BinaryPromiseProblem, FABL.HasPolytimeDistinguisher, FABL.PolytimeMap.comp, FABL.PromiseKarpReducible}
\difficulty{4}
Let the source and target be binary promise problems. If the source reduces to the
target by a polynomial-time map preserving promised yes-instances and no-instances, and
the target has a polynomial-time distinguisher that outputs true on every promised
yes-instance and false on every promised no-instance, then the source has such a
distinguisher as well.
\end{theorem}

\begin{definition}[Encoded maximization problem]
\lean{FABL.EncodedMaximizationProblem}
\statementsource{Analysis-of-Boolean-Functions-by-Ryan-ODonnell}{
  pdf-ffbe9c7b3306-p185-b004,
  pdf-ffbe9c7b3306-p186-b001,
  pdf-ffbe9c7b3306-p186-b003
}
\label{FABL.EncodedMaximizationProblem}
\leanok
An encoded maximization problem consists of a language of well-formed binary instances,
a feasibility relation between binary instances and binary solutions, a rational-valued
objective function on instance–solution pairs, and a rational optimum value for each
binary instance. For every well-formed instance, each feasible solution has objective
value at most the stated optimum.
\end{definition}

\begin{definition}[Threshold approximation guarantee]
\lean{FABL.PolytimeMap.IsThresholdApproximation}
\statementsource{Analysis-of-Boolean-Functions-by-Ryan-ODonnell}{pdf-ffbe9c7b3306-p188-b004}
\label{FABL.PolytimeMap.IsThresholdApproximation}
\leanok
\uses{FABL.EncodedMaximizationProblem, FABL.PolytimeMap}
Let an algorithm be a polynomial-time map on binary strings, and let an encoded
maximization problem have instance language $L$, feasibility relation $F$, objective
value $v$, and optimum value $\mathrm{OPT}$. For rational thresholds $a$ and $b$, the
algorithm is an $(a,b)$ threshold approximation if, for every $x\in L$, its output is
feasible for $x$, and whenever $b\le \mathrm{OPT}(x)$, that output has value at least
$a$.
\end{definition}

\begin{definition}[Polynomial-Time Value Threshold Test]
\lean{FABL.PolytimeMap.ValueThresholdTest}
\statementsource{Analysis-of-Boolean-Functions-by-Ryan-ODonnell}{pdf-ffbe9c7b3306-p196-b008}
\label{FABL.PolytimeMap.ValueThresholdTest}
\leanok
\uses{FABL.EncodedMaximizationProblem, FABL.PolytimeMap}
Let a polynomial-time map compute an output string from each binary input, and let an
encoded maximization problem assign a rational value to each input and output string.
For any rational threshold, a polynomial-time map is a value-threshold test if, on every
input, it outputs the one-bit string containing true exactly when the value of the
computed output is at least the threshold, and false otherwise.
\end{definition}

\begin{theorem}[Weakening approximation thresholds]
\lean{FABL.PolytimeMap.IsThresholdApproximation.weaken}
\label{FABL.PolytimeMap.IsThresholdApproximation.weaken}
\leanok
\uses{FABL.EncodedMaximizationProblem, FABL.PolytimeMap, FABL.PolytimeMap.IsThresholdApproximation}
\difficulty{4}
Let $A$ be a polynomial-time map for an encoded maximization problem, and suppose that
on every well-formed instance its output is feasible and, whenever the optimum is at
least $q$, its output value is at least $r$. If $r'\le r$ and $q\le q'$, then on every
well-formed instance the output is feasible and, whenever the optimum is at least $q'$,
its output value is at least $r'$.
\end{theorem}

\begin{theorem}[Feasible value is bounded by the optimum]
\lean{FABL.PolytimeMap.IsThresholdApproximation.value_le_optimum}
\label{FABL.PolytimeMap.IsThresholdApproximation.value_le_optimum}
\leanok
\uses{FABL.EncodedMaximizationProblem, FABL.PolytimeMap, FABL.PolytimeMap.IsThresholdApproximation}
\difficulty{1}
Let an encoded maximization problem have an algorithm that, on every valid instance,
outputs a feasible solution and, whenever the instance's optimum is at least the
specified optimum threshold, outputs a solution whose value is at least the specified
output threshold. Then, for every valid instance, the algorithm's solution value is at
most the instance's optimum.
\end{theorem}

\begin{definition}[Rational gap promise]
\lean{FABL.EncodedMaximizationProblem.gapPromise}
\label{FABL.EncodedMaximizationProblem.gapPromise}
\leanok
\uses{FABL.BinaryPromiseProblem, FABL.EncodedMaximizationProblem}
Let an encoded maximization problem have a set of well-formed Boolean-string instances
and an optimum value for each Boolean string, and let $s,c\in\mathbb{Q}$ satisfy $s<c$.
Its associated binary promise problem designates as yes-instances the well-formed inputs
whose optimum is at least $c$, and as no-instances the well-formed inputs whose optimum
is at most $s$.
\end{definition}

\begin{definition}[Real gap promise]
\lean{FABL.EncodedMaximizationProblem.gapPromiseReal}
\label{FABL.EncodedMaximizationProblem.gapPromiseReal}
\leanok
\uses{FABL.BinaryPromiseProblem, FABL.EncodedMaximizationProblem}
Let an encoded maximization problem have a set of well-formed Boolean-string instances
and an optimum value for each Boolean string, and let $s,c\in\mathbb{R}$ satisfy $s<c$.
Its associated binary promise problem designates as yes-instances the well-formed inputs
whose optimum, viewed as a real number, is at least $c$, and as no-instances the
well-formed inputs whose optimum, viewed as a real number, is at most $s$.
\end{definition}

\begin{theorem}[Monotonicity of Gap Promises]
\lean{FABL.EncodedMaximizationProblem.gapPromiseReal_reducible_of_thresholds}
\label{FABL.EncodedMaximizationProblem.gapPromiseReal_reducible_of_thresholds}
\leanok
\uses{FABL.EncodedMaximizationProblem, FABL.EncodedMaximizationProblem.gapPromiseReal, FABL.PolytimeMap.id, FABL.PromiseKarpReducible}
\difficulty{4}
Let $P$ be an encoded maximization problem, and let $s<c$ and $s'<c'$ be real thresholds
with $s\le s'$ and $c'\le c$. The promise problem whose yes instances are valid inputs
with optimum at least $c$ and whose no instances are valid inputs with optimum at most
$s$ is polynomial-time Karp reducible to the corresponding promise problem with
thresholds $s'$ and $c'$.
\end{theorem}

\begin{theorem}[Distinguishes gap]
\lean{FABL.PolytimeMap.IsThresholdApproximation.distinguishesGap}
\statementsource{Analysis-of-Boolean-Functions-by-Ryan-ODonnell}{pdf-ffbe9c7b3306-p196-b008}
\label{FABL.PolytimeMap.IsThresholdApproximation.distinguishesGap}
\leanok
\uses{FABL.EncodedMaximizationProblem, FABL.EncodedMaximizationProblem.gapPromise, FABL.PolytimeMap, FABL.PolytimeMap.DistinguishesPromise, FABL.PolytimeMap.IsThresholdApproximation, FABL.PolytimeMap.ValueThresholdTest}
\difficulty{4}
% dataset-pipeline: stub -- prose pending build_informal_statements.py
Statement pending regeneration.
\end{theorem}
